\section{Original line model norms under the same global native homeomorphism}
\label{sec:projective-native-model-line-metrics}

\begin{theorem}[Every original section value under the global homeomorphism {\cite[\S1.2]{BJ22}}]
\label{thm:projective-native-model-global-section-values}
\lean{EnvelopeProperty.ProjectiveNativeModelLineMetrics.nativeModelHomeomorph_chart_val}
\leanok
\uses{thm:projective-fixed-native-analytic-gluing,thm:projective-fixed-native-analytic-values}
The whole global native homeomorphism, viewed in any actual coordinate chart,
evaluates every original regular section by its exact original native value.
\end{theorem}
\begin{proof}
The global map is literally the original chart homeomorphism on that chart.
Its affine-chart square is the original native affine comparison, which
preserves the value of every original section. Compose these two whole-point
identities and evaluate the arbitrary original chart section.
\end{proof}

\begin{theorem}[Actual normalized model charts contain the original proper center {\cite[\S1.2]{BJ22}}]
\label{thm:projective-native-model-value-one-center}
\lean{EnvelopeProperty.ProjectiveNativeModelLineMetrics.center_mem_of_model_coordinate_eq_one}
\leanok
\uses{thm:projective-native-model-global-section-values,scheme:proper-center-bounded-affine-criterion}
If an original coordinate has value one at the image under the same global
homeomorphism, the actual original proper-model center lies in that original
coordinate chart. The center is distinct from the native kernel support;
its chart membership is derived from boundedness of the whole section ring.
\end{theorem}
\begin{proof}
The nonzero coordinate and exact support comparison first put the native
support in the actual original chart. Every element of the whole model chart
algebra is a genuine homogeneous fraction. Its denominator has value one,
while its numerator has value at most one by actual model normalization.
Thus every model chart element has value at most one. Exact original section
values transfer this bound to every element of the entire original native
affine section ring. Its genuine evaluation map factors through the native
valuation ring, yielding an actual affine valuative lift. Proper uniqueness
identifies this lift with the original proper-model lift; hence its center
lies in the original chart. No center-matching premise is assumed.
\end{proof}

\begin{theorem}[Actual original pulled homogeneous line norms {\cite[\S1.2]{BJ22}}]
\label{thm:projective-native-model-homogeneous-line-norms}
\lean{EnvelopeProperty.ProjectiveNativeModelLineMetrics.modelLineNorm_homogeneousSection_eq_chartValue}
\leanok
\uses{thm:projective-native-model-global-section-values,thm:projective-native-model-value-one-center,lem:fc-original-homogeneous-form-frame-coefficient,thm:fc-original-section-pullback-model-norm}
For every original homogeneous polynomial, the norm of its genuine original
global section of the actual pulled-back O(d), evaluated on the original
native residue-field fibre, equals the value of its whole original chart
coefficient at the image under the same global homeomorphism, whenever that
original chart coordinate has model value one. Both the whole original line
and its whole section are retained; no metric comparison is a premise.
\end{theorem}
\begin{proof}
The preceding theorem puts the actual original proper-model center in the
original chart. In its genuine original O(d) frame, the entire original
global homogeneous section has the actual coefficient $F/X_j^d$. Pull this
whole scalar equation back along the original embedding. The genuine proper
model norm is exactly the original native evaluation of this coefficient.
The first theorem identifies that value with its value in the same fixed
model chart under the global homeomorphism. This proves the actual line-norm
identity, including zero sections. Identification with the fixed model's
entire graded section spaces, finite weighted FS functions, and numerical
PSH transport remain subsequent steps.
\end{proof}
